\begin{workedexample}[Worked Example: Range from echo delay]
Radar range follows from the round-trip time: $R = c\,\Delta t/2$. An echo
returning $\Delta t = 20\ \mu\text{s}$ after transmission puts the target at
\[ R = \frac{(3\times10^8)(20\times10^{-6})}{2} = 3000\ \text{m}. \]
Received power falls as $1/R^4$ (two-way spreading times cross-section), so
doubling range cuts echo power by $12\ \text{dB}$.
\end{workedexample}

\begin{keytakeaways}
\begin{itemize}[leftmargin=1.4em,itemsep=2pt]
\item $R = c\Delta t/2$ from the round-trip delay.
\item The radar equation gives echo power $\propto 1/R^4$; doubling range loses $12\ \text{dB}$.
\item Doppler shift $f_d = 2v/\lambda$ gives radial velocity.
\end{itemize}
\end{keytakeaways}

\begin{exercises}
\item Echo delay $60\ \mu\text{s}$: what is the range?
\item Doubling target range changes echo power by how many dB?
\item Doppler shift for $v = 30\ \text{m/s}$ at $\lambda = 3\ \text{cm}$?
\end{exercises}

\furtherreading{Skolnik~\cite{skolnik} (ch.~1--2).}
